\chapter{1972:Christian Anfinsen and Stanford Moore and William H. Stein}

\label{chap:chemistry1972}

The Nobel Prize in Chemistry for the year 1972 was awarded jointly to Christian Anfinsen and Stanford Moore and William H. Stein.

\textbf{Motivation:}

for his work on ribonuclease, especially concerning the connection between the amino acid sequence and the biologically active conformation

\section{Christian Anfinsen}

\subsection*{Early Life and Education}

Christian Anfinsen was born on 1916-03-26 in Monessen, Pennsylvania, USA.

\subsection*{Career and Major Research}

Anfinsen studied at Swarthmore College and received his doctorate from Harvard University in 1943. He worked at the National Institutes of Health in Bethesda, where he studied the pancreatic enzyme ribonuclease. He showed that the enzyme could be unfolded and then refolded spontaneously to its active form.

\subsection*{Nobel Prize-winning Work}

The work for which Christian Anfinsen was awarded the Nobel Prize in Chemistry in 1972 was described by the Nobel Committee as: "for his work on ribonuclease, especially concerning the connection between the amino acid sequence and the biologically active conformation".

Anfinsen demonstrated that the amino acid sequence of ribonuclease determines its biologically active three-dimensional structure, establishing the thermodynamic principle of protein folding.

\subsection*{Significance and Impact}

The finding that a protein folds spontaneously according to its sequence laid the foundation for protein chemistry and for understanding protein folding and misfolding diseases.

\subsection*{Later Career and Honors}

Anfinsen remained at the National Institutes of Health for most of his career. He died on 1995-05-14 in Randallstown, Maryland, USA.

\subsection*{Legacy}

Anfinsen's principle that sequence determines structure is a cornerstone of molecular biology and bioinformatics.

\subsection*{Selected Publications}

"Principles That Govern the Folding of Protein Chains" (Science, 1973)

\clearpage

\section{Stanford Moore}

\subsection*{Early Life and Education}

Stanford Moore was born on 1913-09-04 in Chicago, Illinois, USA.

\subsection*{Career and Major Research}

Moore studied at Vanderbilt University and received his doctorate from the University of Wisconsin. He worked at the Rockefeller Institute, later Rockefeller University, in New York for his entire career. With William H. Stein he developed methods of ion-exchange chromatography for the quantitative analysis of amino acids.

\subsection*{Nobel Prize-winning Work}

The work for which Stanford Moore was awarded the Nobel Prize in Chemistry in 1972 was described by the Nobel Committee as: "for his work on ribonuclease, especially concerning the connection between the amino acid sequence and the biologically active conformation".

Moore and Stein determined the complete amino acid sequence of ribonuclease, providing the chemical basis for Anfinsen's studies of its folding.

\subsection*{Significance and Impact}

Their chromatographic methods and automatic amino acid analyzers made protein sequence determination practical and transformed biochemistry.

\subsection*{Later Career and Honors}

Moore remained at Rockefeller University until his retirement. He died on 1982-08-23 in New York City, USA.

\subsection*{Legacy}

Moore and Stein's analytical methods are the forerunners of modern protein analysis and proteomics.

\subsection*{Selected Publications}

"Photometric Ninhydrin Method for Use in the Chromatography of Amino Acids" (Journal of Biological Chemistry, 1948)

"Automatic Recording Apparatus for Use in the Chromatography of Amino Acids" (Analytical Chemistry, 1958)

\clearpage

\section{William H. Stein}

\subsection*{Early Life and Education}

William H. Stein was born on 1911-06-25 in New York City, New York, USA.

\subsection*{Career and Major Research}

Stein studied at Harvard College and received his doctorate from Columbia University. He joined the Rockefeller Institute in 1938 and remained there for his entire career. With Stanford Moore he developed chromatographic methods for analyzing amino acids and proteins.

\subsection*{Nobel Prize-winning Work}

The work for which William H. Stein was awarded the Nobel Prize in Chemistry in 1972 was described by the Nobel Committee as: "for his work on ribonuclease, especially concerning the connection between the amino acid sequence and the biologically active conformation".

Stein and Moore established the amino acid sequence of ribonuclease and related the structure of the enzyme to its catalytic activity.

\subsection*{Significance and Impact}

Their determination of a complete enzyme sequence, achieved through quantitative chromatography, was a landmark in protein chemistry.

\subsection*{Later Career and Honors}

Stein remained at Rockefeller University for the rest of his career. He died on 1980-02-02 in New York City, USA.

\subsection*{Legacy}

Stein's analytical techniques helped establish the modern era of protein science.

\subsection*{Selected Publications}

"Photometric Ninhydrin Method for Use in the Chromatography of Amino Acids" (Journal of Biological Chemistry, 1948)

"Automatic Recording Apparatus for Use in the Chromatography of Amino Acids" (Analytical Chemistry, 1958)

\clearpage

\section*{Chapter Summary}

The 1972 prize recognized Anfinsen, Moore, and Stein for their work on ribonuclease. Their research connected the amino acid sequence of an enzyme to its biologically active three-dimensional structure.

